\documentclass[11pt]{article}
\usepackage[margin=2.5cm]{geometry}
\usepackage{amsmath,amssymb}
\usepackage{hyperref}
\usepackage{booktabs}
\hypersetup{colorlinks=true,linkcolor=blue,urlcolor=blue}
\newcommand{\code}[1]{\texttt{#1}}

\title{$c\to D^0$ Fragmentation Functions}
\author{}
\date{}

\begin{document}
\maketitle

\noindent The code offers three interchangeable models for the charm-quark-to-$D^0$ fragmentation function $D(z)$, selected by \code{par->frag\_type} (\code{FragmentationType::BCFY} / \code{KniehlKramer} / \code{LHAPDF}, default \code{BCFY}). In every case, the cross-section integrand uses the combination
\begin{equation}
  \text{weight}(z_h) = \frac{D(z_h)}{z_h^2}, \qquad z_h\in[z_{\min},z_{\max}]=[0.05,\,1]
  \label{eq:weight}
\end{equation}
\hfill{\footnotesize\code{integrand.cpp:31--43}, $z_{\min},z_{\max}$: \code{main.cpp:96--97}}
where $z_h$ is the momentum fraction of the charm quark carried by the $D^0$ meson ($p_{D^0}=z_h\,p_c$).

\section{BCFY (Braaten--Cheung--Fleming--Yuan) -- default}
This model splits the fragmenting charm quark into a pseudoscalar piece ($D^0$ directly) and a vector piece ($D^{*0}\to D^0$), each following the perturbative-QCD BCFY functional form, built from a single non-perturbative shape parameter $r$.

\paragraph{Pseudoscalar function $D_P(z,r)$.}
\begin{equation}
  D_P(z,r) = \frac{r\,z\,(1-z)^2}{\big[1-(1-r)z\big]^{6}}\;P(z,r), \qquad 0<z<1
\end{equation}
with
\begin{align}
  P(z,r) = {}& 6 - 18(1-2r)\,z \nonumber\\
  & + \big(21-74r+68r^2\big)\,z^2 \nonumber\\
  & - 2(1-r)\big(6-19r+18r^2\big)\,z^3 \nonumber\\
  & + 3(1-r)^2\big(1-2r+2r^2\big)\,z^4
\end{align}
\hfill{\footnotesize\code{fragmentation.cpp:9--23}}

\paragraph{Vector function $D_V(z,r)$.}
\begin{equation}
  D_V(z,r) = \frac{3\,r\,z\,(1-z)^2}{\big[1-(1-r)z\big]^{6}}\;V(z,r), \qquad 0<z<1
\end{equation}
with
\begin{align}
  V(z,r) = {}& 2 - 2(3-2r)\,z \nonumber\\
  & + 3\big(3-2r+4r^2\big)\,z^2 \nonumber\\
  & - 2(1-r)\big(4-r+2r^2\big)\,z^3 \nonumber\\
  & + (1-r)^2\big(3-2r+2r^2\big)\,z^4
\end{align}
\hfill{\footnotesize\code{fragmentation.cpp:28--41}}

\paragraph{Combination into $D_c^{D^0}(z)$.}
The two channels are combined with branching-fraction-like weights, with the vector piece rescaled in argument by the $D^{*0}/D^0$ mass ratio (to convert "fraction of $D^{*0}$ momentum" into "fraction of $D^0$ momentum") and restricted to the kinematic threshold $z\le m_{D^0}/m_{D^{*0}}$:
\begin{equation}
  \boxed{\;D_c^{D^0}(z) = 0.168\,D_P(z,r) \;+\; 0.39\,\theta\!\left(\frac{m_{D^0}}{m_{D^{*0}}}-z\right)\frac{m_{D^{*0}}}{m_{D^0}}\,D_V\!\left(\frac{m_{D^{*0}}}{m_{D^0}}\,z,\,r\right)\;}
\end{equation}
\hfill{\footnotesize\code{fragmentation.cpp:45--63}}
\begin{center}
\begin{tabular}{@{}ll@{}}
\toprule
Parameter & Value \\
\midrule
$r$ (non-perturbative shape) & $0.1$ \quad\code{main.cpp:93} \\
$m_{D^0}$ & $1.8648$~GeV \\
$m_{D^{*0}}$ & $2.0067$~GeV \\
Pseudoscalar weight & $0.168$ \\
Vector weight & $0.39$ \\
\bottomrule
\end{tabular}
\end{center}

\section{Kniehl--Kramer}
A simpler, purely phenomenological two-parameter form:
\begin{equation}
  \boxed{\;D(z) = N\,\frac{z\,(1-z)^2}{\big[(1-z)^2+\epsilon z\big]^2}\;}, \qquad 0<z<1
\end{equation}
\hfill{\footnotesize\code{fragmentation.cpp:66--72}}
\begin{center}
\begin{tabular}{@{}ll@{}}
\toprule
Parameter & Value \\
\midrule
$N$ & $0.694$ \quad\code{main.cpp:94} \\
$\epsilon$ & $0.101$ \quad\code{main.cpp:95} \\
\bottomrule
\end{tabular}
\end{center}

\section{LHAPDF}
Rather than an analytic form, $D(z)$ is read from an external LHAPDF fragmentation-function grid (member file, default \code{data/prompt-D0-1-109/prompt-D0-1-109\_0000.dat}) for the charm-quark flavor (PDG id $4$), interpolated in $z$ at a fixed factorization scale
\begin{equation}
  Q = \kappa\,\sqrt{p_{D^0}^2+m_c^2}, \qquad \kappa=\texttt{SCALE\_FACTOR}\ (\text{default }1)
\end{equation}
with a floor $Q\ge m_c$ (below the charm-production threshold the grid is undefined, so the scale is not allowed to drop under $m_c$ even for $\kappa=0.5$ at very low $p_{D^0}$). $\kappa=0.5,\,2$ give the conventional scale-variation envelope around the central ($\kappa=1$) result; the 101 members of the grid (one central + 100 replicas) give the PDF-style replica uncertainty band when swapped in one at a time.
\hfill{\footnotesize\code{main.cpp:107--128}}

\end{document}
